% !TeX program = pdflatex
\documentclass[12pt,a4paper]{article}

\usepackage[T2A]{fontenc}
\usepackage[utf8]{inputenc}
\usepackage[english,russian]{babel}
\usepackage{amsmath,amssymb}
\usepackage{graphicx}
\usepackage{hyperref}
\usepackage{geometry}
\geometry{left=2.5cm,right=1.5cm,top=2cm,bottom=2cm}
\hypersetup{colorlinks=true,linkcolor=black,citecolor=black,urlcolor=blue}

\graphicspath{{figures/}}

\newcommand{\avg}{\overline}

\title{\large\bfseries Воспроизводимая оценка детекции и трекинга яблок\\ для подсчёта плодов по видео пролёта БПЛА}
\author{}
\date{}

\begin{document}
\maketitle

\begin{abstract}
Рассматривается задача автоматического подсчёта плодов яблони по
видеозаписи пролёта беспилотного летательного аппарата вдоль ряда деревьев.
Предложен воспроизводимый программный конвейер, разделяющий детекцию,
многокадровое отслеживание и подсчёт единиц результата. Для честного
сравнения трекеров детекции кэшируются и повторно используются всеми
методами; сырые треки оцениваются стандартными MOT-метриками, а объединение
фрагментов после трекинга в основном сравнении не выполняется. Для датасета
VideoFruitCounting единицей счёта является stretch ряда, а не отдельное
дерево; используется общий алгоритм назначения по кадрам маркеров границ.
Оценка абсолютной ошибки счёта планируется на вручную размеченном
контрольном подмножестве, поскольку SFRAME не содержит чисел плодов.
Численные результаты в статье должны заполняться только после завершения
этих контрольных экспериментов; неподтверждённые результаты предварительных
прогонов в текст не включаются.
\end{abstract}

\noindent\textbf{Ключевые слова:} компьютерное зрение, сельское хозяйство,
подсчёт плодов, многокадровое отслеживание, беспилотный летательный
аппарат, воспроизводимость.

\section{Введение}

Оценка числа плодов в плодовом саду необходима для планирования сбора,
логистики и управления агротехническими операциями. Ручной подсчёт
трудоёмок, а фотосъёмка с беспилотного летательного аппарата (БПЛА)
позволяет получить серию изображений всего ряда без остановки движения.

Главная трудность видеоподсчёта состоит в том, что один физический плод
наблюдается на многих кадрах и при этом смещается из-за движения камеры.
Простое суммирование детекций по кадрам многократно завышает результат.
Поэтому требуется трекинг, сохраняющий идентификаторы объектов между
кадрами, и формальное правило перехода от набора треков к итоговому счёту.
При этом число созданных треков само по себе не равно числу плодов: один
плод может давать несколько фрагментов, а ошибочный трек может объединять
разные плоды.

Цель работы --- построить и воспроизводимо оценить конвейер подсчёта
яблок по видео пролёта, включающий: (1) нейросетевую и цветовую детекцию;
(2) сравнение реализаций ByteTrack, BoT-SORT и DeepSORT на одинаковых
детекциях; (3) формальную привязку результатов к единицам ряда;
(4) стандартные метрики детекции, трекинга и счёта; (5) механизм полной
фиксации версий, параметров и входных данных прогона.

В отличие от первоначального прототипа, в настоящей версии не смешиваются
разнородные датасеты, не используются непостоянные идентификаторы COCO как
треки и не публикуются метрики без соответствующего эталона. Прогноз массы
урожая рассматривается отдельно и считается валидированным только при
наличии независимого полевого контроля числа плодов и массы.

\section{Постановка задачи и обзор}

На входе --- последовательность кадров пролёта вдоль ряда яблонь. На выходе
требуется число уникальных плодов по каждой выбранной единице ряда и
сопровождающие метрики качества. Под уникальным плодом понимается
результат корректного трека, а не гарантированная физическая идентичность:
окончательная оценка точности определяется эталонной разметкой.

Базовый подход VideoFruitCounting использует детектор YOLO и ByteTrack, а
результаты привязывает к stretch'ам ряда по кадрам, в которых центр кадра
пересекает маркер границы \cite{felip2026videofruit}. ByteTrack ассоциирует
как высоко-, так и низкоуверенные детекции \cite{zhang2021bytetrack}.
BoT-SORT дополнительно учитывает движение камеры и внешний вид объектов
\cite{aharon2022botsort}, а DeepSORT применяет глубинные признаки для
ассоциации гипотез \cite{wojke2017deepsort}. Сравнение таких схем для
яблоневых садов выполнялось в работе Villacrés et al.
\cite{villacres2023apple}. Отдельная задача --- устойчивое отслеживание
при низкой частоте кадров и периодических окклюзиях, рассмотренная Matos
et al. \cite{matos2024tracking}; в настоящем этапе эта компенсация выделена
как последующая модификация и не подменяет оценку базовых трекеров.

Датасет MinneApple \cite{haeni2019minneapple} предоставляет статические
изображения с разметкой плодов и пригоден для обучения и валидации
детектора. Его идентификаторы аннотаций COCO не постоянны во времени и
поэтому не могут служить эталоном трекинга. Сборка ``видео'' из train-кадров
не используется для оценки обобщения. T1 (Zenodo 7384297) --- отдельный
MOT-набор и не относится к VideoFruitCounting. VideoFruitCounting содержит
16 видео, две камеры и stretch'ы ряда; файл SFRAME задаёт 22 маркера и 21
stretch, но не содержит эталонных чисел плодов.

\section{Метод}

\subsection{Архитектура}

Конвейер состоит из следующих компонентов:
\begin{enumerate}
  \item единый источник кадров с фиксированным порядком, шагом и поворотом;
  \item детектор HSV или YOLO, результат которого сохраняется в общий
        сжатый кэш для всех трекеров;
  \item один из внешних трекеров ByteTrack, BoT-SORT или DeepSORT;
  \item сохранение сырых пофреймовых треков до любой постобработки;
  \item фильтрация треков по фактически выданным наблюдениям и подсчёт;
  \item назначение единиц счёта: stretch по SFRAME либо дерево по отдельной
        карте с корректным эталоном;
  \item расчёт метрик и запись полного манифеста воспроизводимости.
\end{enumerate}

Общий кэш исключает ситуацию, при которой разница трекеров частично
обусловлена повторной детекцией. Кэш содержит рамки, уверенности, номера
кадров и метаданные конфигурации; загрузка проверяет совпадение этих
метаданных. Каждый прогон сохраняет версии пакетов, Git commit, SHA-256
видео и кэша, разрешённую конфигурацию и статус манифеста.

\subsection{Детекция}

Цветовой HSV-детектор служит простым базовым методом и позволяет проверить
весь путь обработки без GPU. Основные эксперименты выполняются двумя
обученными вариантами YOLO: компактной моделью на входе $640\times640$ и
более сильной моделью на входе $1280\times1280$. Для MinneApple обучение и
оценка разделены на train/val; отдельные веса не переносятся на другие
датасеты без явного указания в конфигурации.

\subsection{Трекинг и счёт без повторов}

Локальное сравнение данного этапа включает реализации BoT-SORT и DeepSORT,
получающие идентичные детекции. ByteTrack также подключён к общему
интерфейсу и кэшу, но выделен партнёрскому этапу, чтобы не дублировать
прогоны и не смешивать ответственность за эксперимент. Трек учитывается
после первого фактически выданного наблюдения ($min\_obs=1$), поскольку
при min\_obs$>1$ разные трекеры по-разному подтверждают гипотезы и сравнение
перестаёт отражать только качество ассоциации. Постобъединение фрагментов
по расстоянию отключено ($merge\_dist=0$): оно может повысить итоговый
счёт, но вносит дополнительную эвристику и должно оцениваться отдельным
ablation-этапом.

Важно различать событие счёта и физическое повторение плода. Метрика
$n\_double\_count\_events$ отражает только повторное событие внутреннего
идентификатора и не доказывает повторный учёт разных физических плодов.
Честная оценка последнего требует эталонных траекторий или вручную
проверенного счёта.

\subsection{Единицы счёта и VideoFruitCounting}

Для VideoFruitCounting stretch не считается отдельным деревом: в одном
участке может находиться группа деревьев. Автоматическое дробление stretch
по размеру рамок не является обнаружением отдельных деревьев и в основном
сравнении отключено. Назначение выполняется воспроизведением reference-алго\-ритма VFC по кадрам маркеров, положению центра рамки и скорости
трека. Результат фиксируется как счёт по 21 stretch. Режим median-frame
сохранён только как диагностическая альтернатива и не смешивается с
основным протоколом.

\subsection{Прогноз массы}

Без полевой калибровки физическая оценка
\[
  \hat M = \frac{N\,\avg m}{\kappa}
\]
является иллюстративной: $N$ --- число учтённых плодов, $\avg m$ ---
средняя масса, $\kappa$ --- предполагаемая доля видимых плодов. Регрессия
массы по числу плодов допустима только на независимом полевом контроле;
качество обобщения оценивается leave-one-out или отдельным holdout-набором.
Поэтому в основной матрице трекеров прогноз урожая отключён.

\section{Экспериментальный протокол}

\subsection{Разделение ролей данных}

\begin{itemize}
  \item MinneApple val --- только метрики детекции;
  \item T1 --- стандартные MOT-метрики на постоянных идентификаторах;
  \item VideoFruitCounting --- счёт по stretch, абсолютная точность ---
        только на ручном count GT контрольного подмножества;
  \item синтетические данные --- регрессионные тесты и диагностика, не
        заменяют полевую валидацию.
\end{itemize}

\subsection{Матрица}

Для каждого источника выполняются комбинации трёх детекторов (HSV,
YOLO-модель 640, YOLO-модель 1280) и двух локальных трекеров BoT-SORT и
DeepSORT. ByteTrack подключён через тот же интерфейс и может быть добавлен
в матрицу одной строкой без изменения формата результатов. Для одного видео
и детектора кэш создаётся один раз; все трекеры получают его без изменения.
Параметры настраиваются только на validation-подмножестве и не подгоняются
под контрольный count GT.

\begin{table}[htbp]
\centering
\caption{План локального сравнения методов}
\label{tab:plan}
\begin{tabular}{lll}
\hline
Источник & Детекция & Трекинг \\
\hline
MinneApple val & HSV, YOLO640, YOLO1280 & не оценивается \\
T1 & один общий кэш на детектор & BoT-SORT, DeepSORT \\
VFC & один общий кэш на видео и детектор & BoT-SORT, DeepSORT \\
\hline
\end{tabular}
\end{table}

\subsection{Метрики}

Детекция оценивается AP@0.50, precision, recall, F1 и ошибкой числа
детекций. Трекинг --- MOTA, IDF1 и IDSW через MOT-совместимую библиотеку.
Счёт --- MAE, RMSE, смещение, MAPE и относительная ошибка итогового счёта
по единице и видео. Производительность разделяется на время детекции и
время трекинга. Результаты комбинаций без валидного манифеста и точного
совпадения конфигурации исключаются из сводной таблицы.

\section{Результаты}

Численные результаты в этот раздел вносятся только после завершения новых
прогонов по описанному протоколу. Ниже приведены обязательные таблицы,
которые должны заполняться из `matrix\_results.csv` и сохранённых summary,
без переноса результатов прежней реализации.

\begin{table}[htbp]
\centering
\caption{Детекция на MinneApple val (заполняется после прогона)}
\label{tab:detection}
\begin{tabular}{lccccc}
\hline
Детектор & AP@0.50 & Precision & Recall & F1 & count-MAE \\
\hline
HSV & --- & --- & --- & --- & --- \\
YOLO @640 & --- & --- & --- & --- & --- \\
YOLO @1280 & --- & --- & --- & --- & --- \\
\hline
\end{tabular}
\end{table}

\begin{table}[htbp]
\centering
\caption{Трекинг на T1 при общих детекциях (заполняется после прогона)}
\label{tab:mot}
\begin{tabular}{llccc}
\hline
Детектор & Трекер & MOTA & IDF1 & IDSW \\
\hline
\multicolumn{5}{l}{заполняется автоматически из валидных прогонов}\\
\hline
\end{tabular}
\end{table}

\begin{table}[htbp]
\centering
\caption{Счёт по stretch на VFC (только при наличии ручного count GT)}
\label{tab:count}
\begin{tabular}{llcccc}
\hline
Видео & Комбинация & MAE & RMSE & Bias & Ошибка итога, \% \\
\hline
\multicolumn{6}{l}{не заполняется до подготовки контрольной разметки}\\
\hline
\end{tabular}
\end{table}

\section{Обсуждение и ограничения}

Даже высокие MOTA и IDF1 не гарантируют правильный итоговый счёт: ошибки
детекции, фрагментация и слияние треков влияют на число учтённых объектов
по-разному. Поэтому выбор лучшей комбинации по косвенным признакам,
например только по числу созданных треков, недопустим. Для VFC выбор
наилучшей пары детектор--трекер должен опираться на вручную размеченный
счёт контрольных stretch'ей.

Эвристическое разделение stretch на ``деревья'' по размеру рамок не имеет
полевого эталона и не используется для выводов о числе деревьев. Прогноз
массы без взвешивания контрольных деревьев носит демонстрационный характер.
Компенсация движения камеры и восстановление идентификаторов после длительных
пропусков соответствуют отдельному третьему этапу и не должны приписываться
базовым прогонам ByteTrack, BoT-SORT и DeepSORT.

\section{Заключение}

Сформирован воспроизводимый конвейер и экспериментальный протокол для
сравнения детекторов и трекеров при подсчёте яблок по видео. Ключевые
исправления: строгое разделение ролей MinneApple, T1 и VideoFruitCounting;
общий кэш детекций; внешние реализации трекеров; сырые MOT-метрики без
необоснованного слияния; корректная трактовка stretch; отказ от публикации
прогноза урожая без полевой калибровки; полная фиксация условий прогона.
Дальнейшие работы включают ручную count-разметку контрольного подмножества
VFC, полные матрицы прогонов, ablation min\_obs и merge, а затем третий
этап компенсации движения камеры и восстановления идентификаторов.

\begin{thebibliography}{99}
\bibitem{haeni2019minneapple}
H{\"a}ni N., Roy P., Isler V. MinneApple: a benchmark dataset for apple
detection and segmentation. arXiv:1909.06441, 2019.

\bibitem{felip2026videofruit}
Felip-Pom{\'e}s M., Gen{\'e}-Mola J., Arn{\'o} J., Lordan J.,
Ruiz-Hidalgo J., Morros J.-R., Net Barn{\'e}s F., Gregorio E.
Video-based fruit detection and tracking: effects of scanning conditions on
counting and mapping accuracy. VideoFruitCounting, 2026.
\url{https://github.com/imatge-upc/VideoFruitCounting}

\bibitem{zhang2021bytetrack}
Zhang Y., Sun P., Jiang Y., Yu D., Weng F., Yuan Z., Luo P., Liu W., Wang X.
ByteTrack: multi-object tracking by associating every detection box.
arXiv:2110.06864, 2021.

\bibitem{aharon2022botsort}
Aharon N., Orfaig R., Bobrovsky B.-Z. BoT-SORT: robust associations
multi-pedestrian tracking. arXiv:2206.14651, 2022.

\bibitem{wojke2017deepsort}
Wojke N., Bewley A., Paulus D. Simple online and realtime tracking with a
deep association metric. Proc. IEEE ICIP, 2017, pp. 3645--3649.

\bibitem{villacres2023apple}
Villacr{\'e}s S. et al. Apple orchard production estimation using deep
learning strategies: a comparison of tracking-by-detection algorithms.
Computers and Electronics in Agriculture, 2023, vol. 204, 107513.
\url{https://doi.org/10.1016/j.compag.2022.107513}

\bibitem{matos2024tracking}
Matos B. et al. Tracking and counting apples in orchards under intermittent
occlusions and low frame rates. Proc. CVPR Workshops, 2024.
\url{https://openaccess.thecvf.com/content/CVPR2024W/Vision4Ag/html/Matos_Tracking_and_Counting_Apples_in_Orchards_Under_Intermittent_Occlusions_and_CVPRW_2024_paper.html}

\bibitem{jocher2023ultralytics}
Jocher G., Chaurasia A., Qiu J. Ultralytics YOLO, 2023.
\url{https://github.com/ultralytics/ultralytics}

\bibitem{minneapple2019data}
MinneApple dataset. University of Minnesota, 2019.
\url{https://conservancy.umn.edu/handle/11299/206575}
\end{thebibliography}

\end{document}
